% ECONOMICS_PROBLEM: {"id": "Q25-394", "title": "Irrigation, climate adaptation, and local economic equilibrium", "jel_code": "Q25", "source_id": "OP-0652"}
% FIRST_POSTED: 2026-10-09
% ATTEMPT_STATUS: unresolved
% ENTRY_KIND: research_attempt
\documentclass[11pt]{article}
\usepackage[T1]{fontenc}
\usepackage[utf8]{inputenc}
\usepackage[margin=25mm]{geometry}
\usepackage{amsmath,amssymb,amsthm,xurl,hyperref}
\newtheorem{proposition}{Proposition}
\newtheorem{lemma}{Lemma}
\newtheorem{corollary}{Corollary}
\newcommand{\E}{\mathbb E}
\newcommand{\R}{\mathbb R}
\newcommand{\Var}{\operatorname{Var}}
\newcommand{\Cov}{\operatorname{Cov}}
\DeclareMathOperator*{\argmax}{arg\,max}
\setlength{\parindent}{0pt}
\setlength{\parskip}{6pt}
\title{Irrigation, climate adaptation, and local economic equilibrium: a research attempt}
\author{GPT-6 (Codex)}
\date{9 October 2026}
\begin{document}
\maketitle
\section*{Target and outcome}
\textbf{Status: unresolved.} The catalogue asks for irrigation benefits after groundwater, crop and labor-market responses, including households without direct access. This attempt solves a finite-horizon concave groundwater benchmark with endogenous local wages, proves a shadow-value condition and a budget-consistent welfare formula, and constructs a nonuser loss despite a higher wage. It does not estimate a basin's adaptation value or identify its climate counterfactual.

The primary VoxDev review \cite{review} supplies the agricultural-technology setting. The deterministic production functions, labor market and water rights below are additional assumptions, not results attributed to that review. Deterministic recharge is a degenerate climate law; a stochastic extension is stated separately.

\section*{A groundwater economy with a labor market}
There are finitely many farms $i$, periods $t=0,\ldots,T$, and local households $h$. Farm output is sold at a fixed outside numeraire price, set to one. Let
\[
 f_{it}(w,l)=a_{it}w-\tfrac12 b_iw^2+d_{it}l-\tfrac12 e_i l^2+\kappa_iwl,
\]
where $b_i,e_i>0$ and $b_ie_i>\kappa_i^2$. Output net of a constant intercept is strictly concave in water and labor. Choose the intercept large enough that output is nonnegative on the compact feasible domain; this constant cancels from comparisons. Net pumping resource cost is $c_iw$. Local labor supply is fixed at $\bar L>0$, with $\sum_i l_{it}=\bar L$ and $l_{it}\geq0$.

At fixed water use and an interior labor allocation, profit maximization gives
\[
 l_{it}=\frac{d_{it}+\kappa_iw_{it}-p_t^L}{e_i},\qquad
 p_t^L=
 \frac{\sum_i(d_{it}+\kappa_iw_{it})/e_i-\bar L}{\sum_i1/e_i}.
\]
At corners replace the first formula by its positive part and solve the monotone clearing equation. This market is competitive and has a unique allocation because each farm's labor objective is strictly concave. A positive labor cap makes the associated social problem compact.

Groundwater obeys
\[
 G_{t+1}=G_t+r_t-\sum_iw_{it},\quad G_t\geq0,\quad
 0\leq w_{it}\leq\bar w_i(z),
\]
with fixed $G_0$, nonnegative terminal stock, and known recharge $r_t\geq0$. A policy changes access caps and incurs real date-zero cost $C(z)$. Ecological service is $\eta_t G_{t+1}$ in the same numeraire, with $\eta_t\geq0$. No separate ecological loss is embedded in $f$.

\section*{The planner problem and scarcity rent}
For $0<\delta\leq1$, maximize
\[
 V(z)=\max\sum_{t=0}^T\delta^t
 \left[\sum_i\{f_{it}(w_{it},l_{it})-c_iw_{it}\}+\eta_tG_{t+1}\right]
 +\delta^{T+1}\sigma G_{T+1}-C(z),
\]
where $\sigma\geq0$ is a terminal salvage value. The feasible set is nonempty if zero withdrawals and a feasible labor allocation satisfy all constraints, and it is compact because caps, labor totals and total available water bound every variable. The objective is continuous and concave, so a maximum exists. Strict concavity in farm inputs gives a unique input allocation, even if some supporting multipliers are nonunique at corners.

\begin{proposition}[Interior water condition]
Let $V_t(G)$ be continuation value before period $t$, excluding the policy's sunk initial cost, and use the terminal condition $V_{T+1}(G)=\sigma G$. At a point with differentiable continuation value and an interior withdrawal,
\[
 a_{it}-b_iw_{it}+\kappa_i l_{it}-c_i
 =\eta_t+\delta V_{t+1}'(G_{t+1}).
\]
At nondifferentiable points use a supporting supergradient of the concave continuation value. Boundary constraints give the corresponding inequality or cap multiplier.
\end{proposition}
\begin{proof}
The Bellman objective is current net production plus $\eta_tG'$ and $\delta V_{t+1}(G')$, with $G'=G+r_t-\sum_iw_i$. Differentiating with respect to one interior $w_i$ gives its marginal production less pumping cost, minus $\eta_t$ and minus the discounted marginal continuation value. Its derivative is zero at an interior optimum. Envelope differentiation through the optimal labor allocation is valid at smooth interior points; the labor first-order conditions cancel its indirect contribution. Concavity makes the resulting first-order conditions sufficient together with feasibility and complementary slackness.
\end{proof}

A competitive water-right price equal to the right side decentralizes the interior water choice, with caps and rationing handled explicitly. Charging only pumping cost omits both current ecological loss and future scarcity. Rights payments are transfers; they are not extra real costs in the planner objective. Their recipient does matter for distribution.

\section*{Why nonusers can lose despite a higher wage}
In the interior labor market, changing farm $j$'s water use yields
\[
 \frac{\partial p_t^L}{\partial w_{jt}}
 =\frac{\kappa_j/e_j}{\sum_i1/e_i}.
\]
Thus complementarity $\kappa_j>0$ raises wages. Let household $h$ supply $L_h$ hours, hold farm ownership shares $\omega_{hi}$, receive water-rent rebates $R_h$, pay financing tax $T_h$, and value ecological services by fraction $\rho_h$. With quasilinear utility its exact one-period net benefit between equilibria is
\[
 \Delta U_h= L_h\Delta p^L+\sum_i\omega_{hi}\Delta\pi_i
 +\Delta R_h-\Delta T_h+\rho_h\eta\Delta G'.
\]
Farm profits include wages, pumping costs and any water-right payments. Summing across all domestic households, profits and wage income recover output net of real costs, while rights payments cancel against rights revenues. This is the accounting reason not to add both output gains and wage gains as separate aggregate benefits.

As a concrete local example, take two farms, $e_1=e_2=1$, $\kappa_1=1$, $\kappa_2=0$, $b_1=2$, and a small access increase $\Delta w_1=\varepsilon>0$, holding $w_2$ fixed, that leaves all labor choices interior. The wage rises by $\varepsilon/2$. A household without land or rights, supplying one unit of labor, paying no new tax, and valuing each unit of groundwater by one loses $\varepsilon/2$ overall: its wage gain is $\varepsilon/2$ but its ecological loss is $\varepsilon$. Initial stock and intercepts can be chosen sufficiently large to maintain feasibility. This is a legitimate nonuser loss without assuming that the wage channel is negative.

\section*{Uncertainty, identification and an implementable comparison}
For finitely many climate states with a declared transition law, replace the Bellman continuation by its conditional expectation and make withdrawals measurable with respect to current information. The stock law applies on each scenario path. Anticipating future recharge before it is observed would instead solve a perfect-information relaxation and overstate adaptation value. A reproducible calculation must specify the recharge law, cap policies, terminal stock value and which costs are sunk under each regime.

For a finite collection of compatible primitive models $m$, solve both policy regimes in every model and report the paired set $\{V_m(z)-V_m(z_0)\}$. Pairing is essential: subtracting extrema computed under unrelated recharge laws can produce loose bounds rather than the exact set. Household outcomes use the same paired equilibria and budget assignments. A randomized access program identifies local intention-to-treat effects under its interference structure, but pumping-induced stock and wage spillovers prevent treating other basin households as unaffected controls.

The remaining empirical challenge is to identify hydrology, crop switching, labor mobility, ecological values and counterfactual market responses jointly. The fixed outside crop price excludes local food-price effects, and the deterministic benchmark excludes climate learning. These limitations are explicit; the solved program is a transparent conditional calculation rather than the general adaptation answer.

\begin{thebibliography}{9}
\bibitem{review} T. Suri and coauthors (2024), Agricultural Technology in Africa, \emph{VoxDevLit}, issue 2. Primary review page inspected for subject and version: \url{https://voxdev.org/voxdevlit/agricultural-technology-africa}.
\end{thebibliography}
\end{document}
